\section{Week 6: Bernstein--Vazirani and Simon}\label{sec:week6}

Lectures of September 29 and October 1.

\begin{opentable}
    \openrow{lem:signsum-0929}{Sum of signs $\sum_x (-1)^{x \cdot y}$}
    \openrow{thm:bv-0929}{Bernstein--Vazirani}
    \openrow{thm:simon-run-0929}{One run of Simon's circuit}
    \openrow{thm:simon-0929}{Simon's algorithm}
    \openrow{ex:simon-det-0929}{Exercise: deterministic $\bigO{\sqrt{2^n}}$ algorithm}
    \openrow{thm:simon-rand-ub-1001}{Randomized upper bound}
    \openrow{thm:simon-rand-lb-1001}{Randomized lower bound}
    \openrow{lem:yao-1001}{Yao's min-max principle, easy direction}
    \openrow{lem:collision-1001}{Few queries rarely collide}
    \openrow{thm:simon-det-lb-1001}{Classical query complexity $\bigTh{\sqrt{2^n}}$}
    \openrow{ex:simon-11-1001}{Exercise: Simon's circuit when $f$ is $1$-to-$1$}
    \openrow{thm:indep-1001}{Exercise: random vectors are independent w.p.\ $\geq \sfrac{1}{4}$}
\end{opentable}

\subsection{Running the Algorithms}

\begin{exercisebox}
\begin{pproblem}[Bernstein--Vazirani with $n = 3$]\label{pp:bv}
    Let $f(x) = s \cdot x \bmod 2$ with the hidden string $s = 101$, so $f(x) = x_1 \oplus x_3$.
    \begin{enumerate}[label=(\alph*)]
        \item How does a classical algorithm find $s$?
        \item Run the quantum algorithm.
    \end{enumerate}
\end{pproblem}
\tcblower
\begin{wsolution}
    \uses{\scr{thm:bv-0929}; \scr{lem:signsum-0929}}
    (a) Query the unit vectors: $f(100) = 1$, $f(010) = 0$, $f(001) = 1$. The answers are the bits
    of $s$, so $n = 3$ queries. Fewer cannot work, since each answer is one bit and $s$ has $n$ bits.

    (b) After the first Hadamards and the query, the first register is
    \[\frac{1}{\sqrt{8}} \sum_{x} (-1)^{s \cdot x} \ket{x}
        = \frac{1}{\sqrt{8}} \left(\ket{000} - \ket{001} + \ket{010} - \ket{011} - \ket{100} + \ket{101} - \ket{110} + \ket{111}\right).\]
    After the second Hadamards the amplitude of $\ket{y}$ is
    \[\frac{1}{8} \sum_{x} (-1)^{s \cdot x + y \cdot x} = \frac{1}{8} \sum_{x} (-1)^{(s \oplus y) \cdot x}
        = \begin{cases} 1 & \text{if } y = s, \\ 0 & \text{if } y \neq s, \end{cases}\]
    by the sum-of-signs claim (Sep 29): the sum is $2^n$ if $s \oplus y = 000$ and $0$ otherwise.
    So the measurement gives $y = 101 = s$ with probability $1$, after one query.
\end{wsolution}
\end{exercisebox}

\begin{exercisebox}
\begin{pproblem}[Simon with $n = 3$]\label{pp:simon}
    The function $f \colon \set{0,1}^3 \to \set{0,1}^3$ is given by the table below. Run Simon's
    algorithm and find $s$.
    \begin{center}
        \begin{tabular}{c|cccccccc}
            $x$    & 000 & 001 & 010 & 011 & 100 & 101 & 110 & 111 \\ \hline
            $f(x)$ & 101 & 010 & 000 & 110 & 000 & 110 & 101 & 010
        \end{tabular}
    \end{center}
\end{pproblem}
\tcblower
\begin{wsolution}
    \uses{\scr{thm:simon-run-0929}; \scr{thm:simon-0929}; \scr{thm:indep-1001}}
    \leavevmode
    \begin{enumerate}
        \item After the Hadamards and the query, the state is
              $\frac{1}{\sqrt{8}} \sum_{x} \ket{x}\ket{f(x)}$.
        \item Measure the second register. Suppose the outcome is $010$. The inputs with
              $f(x) = 010$ are $001$ and $111$, so the first register is
              $\frac{1}{\sqrt{2}}\left(\ket{001} + \ket{111}\right)$.
        \item Apply $H^{\otimes 3}$. The amplitude of $\ket{y}$ is
              \[\frac{1}{\sqrt{2}} \cdot \frac{1}{\sqrt{8}} \left((-1)^{001 \cdot y} + (-1)^{111 \cdot y}\right)
                  = \frac{1}{4} (-1)^{001 \cdot y} \left(1 + (-1)^{110 \cdot y}\right).\]
        \item This is $\pm\frac{1}{2}$ if $110 \cdot y = 0$ and $0$ if $110 \cdot y = 1$. So $y$ is
              one of $000, 001, 110, 111$, each with probability $\sfrac{1}{4}$. The state is
              $\frac{1}{2}\left(\ket{000} - \ket{001} + \ket{110} - \ket{111}\right)$.
        \item Every outcome satisfies $y_1 \oplus y_2 = 0$. The same holds for any other outcome in
              step 2, because every pair $\set{x, x \oplus s}$ has the same difference.
        \item Run the circuit $n - 1 = 2$ times. Suppose the outcomes are $y = 111$ and $y = 001$.
              They give the equations $s_1 \oplus s_2 \oplus s_3 = 0$ and $s_3 = 0$.
        \item Solve: $s_3 = 0$ and $s_1 = s_2$. The solutions are $000$ and $110$, so the nonzero
              candidate is $s' = 110$.
        \item Check with two classical queries: $f(000) = 101 = f(110)$. So $f$ is $2$-to-$1$ and
              $s = 110$.
    \end{enumerate}

    \emph{What can go wrong.} If the two outcomes are dependent (for example $000$ and $111$, or
    $001$ twice), the system has more than two solutions and we run again. Two outcomes are
    independent with probability $\frac{3}{4} \cdot \frac{1}{2} = \frac{3}{8}$: the first must be
    nonzero ($3$ of $4$), and the second must be outside $\set{000, y_1}$ ($2$ of $4$). This matches
    the bound of at least $\sfrac{1}{4}$ from Oct 1.

    \emph{If $f$ were $1$-to-$1$.} Step 2 would leave a single $\ket{x}$, every $y$ would have
    probability $\sfrac{1}{8}$, and the check in step 8 would fail, so the algorithm would answer
    $s = 000$.
\end{wsolution}
\end{exercisebox}


\begin{exercisebox}
\begin{pproblem}[Simon's circuit for a $1$-to-$1$ $f$, $n = 2$]\label{wp:simon11}
    Let $f \colon \set{0,1}^2 \to \set{0,1}^2$ be $1$-to-$1$. After the second register is measured
    with outcome $f(x_0)$, what is the state, and with what probability is each $y$ measured?
\end{pproblem}
\tcblower
\begin{wsolution}
    \uses{\scr{ex:simon-11-1001}; \scr{lem:wh-0924}}
    Only $x_0$ has the value $f(x_0)$, so the first register is the single basis state
    $\ket{x_0}$. Then
    $H^{\otimes 2}\ket{x_0} = \frac{1}{2} \sum_y (-1)^{x_0 \cdot y} \ket{y}$, so each of the $4$
    strings $y$ has probability $\left(\frac{1}{2}\right)^2 = \frac{1}{4}$, and these add up to
    $1$. For general $n$ each $y$ has probability $\sfrac{1}{2^n}$. Unlike the $2$-to-$1$ case, no
    $y$ is ruled out.
\end{wsolution}
\end{exercisebox}

\begin{exercisebox}
\begin{pproblem}[Independence for $n = 4$]\label{wp:indep4}
    Simon's algorithm with $n = 4$ collects $y_1, y_2, y_3$. Find the probability that they are
    linearly independent if $f$ is $2$-to-$1$, and if $f$ is $1$-to-$1$. How many attempts are
    expected in the $2$-to-$1$ case?
\end{pproblem}
\tcblower
\begin{wsolution}
    \uses{\scr{thm:indep-1001}}
    \emph{$2$-to-$1$.} The $y$'s are uniform in $V = \set{y \suchthat s \cdot y = 0}$, of dimension
    $d = 3$. Part 1 of the exercise gives
    \[\left(1 - \frac{1}{2}\right)\left(1 - \frac{1}{4}\right)\left(1 - \frac{1}{8}\right)
        = \frac{1}{2} \cdot \frac{3}{4} \cdot \frac{7}{8} = \frac{21}{64} \approx 0.33 \geq \frac{1}{4}.\]
    Each factor is the probability that the next vector is outside the span of the earlier ones:
    the span of $i - 1$ independent vectors has $2^{i-1}$ of the $2^3$ elements of $V$.

    \emph{$1$-to-$1$.} Now $V = \set{0,1}^4$, $d = 4$, and we need $d - 1 = 3$ independent
    vectors. Part 2 gives
    $\frac{3}{4} \cdot \frac{7}{8} \cdot \frac{15}{16} = \frac{315}{512} \approx 0.62 \geq \frac{1}{2}$.

    \emph{Attempts.} Each attempt succeeds with probability $\sfrac{21}{64}$, independently, so the
    number of attempts is geometric with mean $\sfrac{64}{21} \approx 3.05$ (\scr{thm:geometric-0827}).
\end{wsolution}
\end{exercisebox}

\begin{exercisebox}
\begin{pproblem}[Deterministic Simon for $n = 4$]\label{pp:det-simon}
    Run the deterministic algorithm of the Sep 29 exercise for $n = 4$.
    \begin{enumerate}[label=(\alph*)]
        \item Which inputs does it query, and how many?
        \item Which two queries collide if $s = 0110$? If $s = 1100$?
        \item How many queries for $n = 20$, against the obvious deterministic algorithm and the
              quantum one?
    \end{enumerate}
\end{pproblem}
\tcblower
\begin{wsolution}
    \uses{\scr{ex:simon-det-0929}}
    The algorithm splits $x$ into its first $m = \ceil{\sfrac{n}{2}}$ bits $u$ and its last $n - m$
    bits $v$, and queries every string of the form $(u, 0 \dots 0)$ and every string of the form
    $(0 \dots 0, v)$.

    (a) Here $m = 2$. The strings $u00$ are $0000, 0100, 1000, 1100$. The strings $00v$ are
    $0000, 0001, 0010, 0011$. The string $0000$ is in both, so there are $4 + 4 - 1 = 7$ queries.

    (b) $s = 0110$ has $u = 01$ and $v = 10$. The queries $0100$ and $0010$ have XOR $0110 = s$, so
    $f(0100) = f(0010)$. $s = 1100$ has $u = 11$ and $v = 00$. The queries $1100$ and $0000$ have
    XOR $1100 = s$, so $f(1100) = f(0000)$.

    In general every $s = (u, v) \neq 0000$ is the XOR of the two queried strings $(u, 00)$ and
    $(00, v)$, and these are different because $s \neq 0000$. So a collision always appears, and by
    the promise any collision $x, x'$ has $x \oplus x' = s$. If no two answers agree, $f$ is
    one-to-one and $s = 0000$.

    (c) The count is $2^m + 2^{n-m} - 1$. For $n = 20$ that is $1024 + 1024 - 1 = 2047$. Querying
    inputs until two answers agree needs $2^{19} + 1 = 524{,}289$ queries in the worst case. Simon's
    algorithm makes $n + 1 = 21$.

    So the deterministic count is about $2\sqrt{N}$, the same order as the randomized lower bound of \cref{pp:simon-lb}, and the quantum count is about $\log N$.
\end{wsolution}
\end{exercisebox}


\subsection{Classical Lower Bounds and the Separation}

\begin{exercisebox}
\begin{pproblem}[The same for Monte Carlo algorithms]\label{pp:yao-mc}
    Show that a randomized search algorithm with success probability at least $\sfrac{2}{3}$ on
    every input needs at least $\sfrac{2n}{3} - 1$ queries.
\end{pproblem}
\tcblower
\begin{wsolution}
    \uses{\scr{lem:yao-1001}; \scr{thm:detsearch-0825}}
    \leavevmode
    \begin{enumerate}
        \item Use the same distribution $\mu$. Take a deterministic algorithm with $T$ queries, in
              the fixed order $q_1, \dots, q_T$ while the answers are ``no''.
        \item It finds $x$ if $x$ is in one of these $T$ slots: probability $\sfrac{T}{n}$.
        \item Otherwise it has seen $T$ answers of ``no'' and outputs one fixed slot. That slot is
              the right one with probability at most $\sfrac{1}{n}$ overall.
        \item So it succeeds with probability at most $\sfrac{(T+1)}{n}$ under $\mu$.
        \item If $T + 1 < \sfrac{2n}{3}$, every deterministic algorithm with $T$ queries errs with
              probability more than $\sfrac{1}{3}$ under $\mu$. By Yao's principle (the form proved
              on Oct 1) every randomized algorithm with $T$ queries errs with probability more than
              $\sfrac{1}{3}$ on some input. \qedhere
    \end{enumerate}
\end{wsolution}
\end{exercisebox}

\begin{exercisebox}
\begin{pproblem}[Simon's problem, in outline]\label{pp:simon-lb}
    State the hard distribution for Simon's problem and the key estimate. Evaluate it for $n = 10$
    and $16$ queries.
\end{pproblem}
\tcblower
\begin{wsolution}
    \uses{\scr{thm:simon-rand-lb-1001}; \scr{lem:collision-1001}}
    \leavevmode
    \begin{enumerate}
        \item Distribution: $s$ is uniform among the $N - 1$ nonzero strings, and $f$ is a uniformly
              random $2$-to-$1$ function with period $s$.
        \item An algorithm learns $s$ only from a collision, two queries $x$ and $x \oplus s$.
        \item With $k$ distinct queries the probability of a collision is at most
              $\sfrac{k(k-1)}{N}$, since there are $\sfrac{k(k-1)}{2}$ pairs and each pair rules in
              one value of $s$ out of at least $\sfrac{N}{2}$.
        \item Without a collision, $s$ is still uniform over at least $\sfrac{N}{2}$ strings, so a
              guess is right with probability at most $\sfrac{2}{N}$.
        \item The success probability is at most $\sfrac{(k(k-1) + 2)}{N}$.
    \end{enumerate}
    For $n = 10$, $N = 1024$ and $k = 16 = \sfrac{\sqrt{N}}{2}$: at most
    $\sfrac{(240 + 2)}{1024} \approx 0.24$. So $\sfrac{\sqrt{N}}{2}$ queries are not enough for
    success probability $\sfrac{2}{3}$, and the classical complexity is $\bigOm{\sqrt{2^n}}$.
\end{wsolution}
\end{exercisebox}

\begin{exercisebox}
\begin{pproblem}[Why Simon is different]\label{pp:simon-sep}
    Explain why Simon's problem gives a real exponential separation.
\end{pproblem}
\tcblower
\begin{wsolution}
    \uses{\scr{thm:simon-rand-lb-1001}; \scr{thm:simon-det-lb-1001}}
    The Oct 1 lower bound is for randomized algorithms with bounded error: success probability
    $\sfrac{2}{3}$ needs more than $\sfrac{\sqrt{N}}{2}$ queries. The quantum algorithm needs
    $\bigO{n}$ queries. So the gap is $\sqrt{2^n}$ against $n$ even when the classical algorithm may
    use randomness and may err. A quantum advantage has to be shown against randomized algorithms,
    because they are the strongest classical model (Aug 25).
\end{wsolution}
\end{exercisebox}



